\section{Exercices d'entraînement corrigés}

\noindent\textit{Essayez de résoudre chaque exercice seul (avec le formulaire) avant de lire la correction.
Les 4 exercices couvrent les 3 types qui reviennent le plus souvent à l'examen : cycle moteur/turbine à gaz,
moteur à piston (diesel, 2 temps), et cycle frigorifique.}

\subsection{Exercice 1 -- Cycle moteur à turbine à gaz (double détente)}

\begin{enonce}
On observe un cycle moteur utilisant de l'air (gaz parfait, $c_p=1004$ J/kg·K, $\gamma=1{,}4$), débit
massique $\dot m=35$ kg/s :
\begin{itemize}[leftmargin=1.4em,itemsep=1pt]
\item 1-2 : compression adiabatique réversible de 1 atm à 4,1 atm ; $T_1=15\,^\circ$C.
\item 2-3 : échauffement isobare jusqu'à $T_3=750\,^\circ$C.
\item 3-4 : détente adiabatique réversible dans une 1\textsuperscript{ère} turbine ; $p_4$ inconnue, mais
tout le travail de cette turbine sert uniquement à entraîner le compresseur (1-2).
\item 4-5 : échauffement isobare jusqu'à $T_5=700\,^\circ$C.
\item 5-6 : détente adiabatique réversible dans une 2\textsuperscript{nde} turbine, jusqu'à $p_6=p_1$.
\item 6-1 : refroidissement isobare (fermeture du cycle).
\end{itemize}
\emph{Trouver : $T_2,T_4,T_6$ et $p_4$ ; le rendement du cycle ; la puissance utile de la 2\textsuperscript{nde}
turbine.}
\end{enonce}

\begin{correction}
\textbf{Point 2} (compression adiabatique, $T_1=288{,}15$ K) :
\[T_2=T_1\left(\frac{p_2}{p_1}\right)^{\frac{\gamma-1}{\gamma}}=288{,}15\times4{,}1^{0{,}2857}
=\mathbf{431{,}2\text{ K}}\]

\textbf{Point 4} : le travail de la turbine 1 égale (en valeur absolue) celui du compresseur, avec le même
débit dans les deux : $c_p(T_3-T_4)=c_p(T_2-T_1)\Rightarrow T_3-T_4=T_2-T_1$.
Avec $T_3=1023{,}15$ K :
\[T_4=T_3-(T_2-T_1)=1023{,}15-(431{,}2-288{,}15)=\mathbf{880{,}1\text{ K}}\]

\textbf{Pression $p_4$} (adiabatique 3-4) : $\dfrac{T_4}{T_3}=\left(\dfrac{p_4}{p_2}\right)^{0{,}2857}$
\[p_4=p_2\left(\frac{T_4}{T_3}\right)^{1/0{,}2857}=4{,}1\times\left(\frac{880{,}1}{1023{,}15}\right)^{3{,}5}
=\mathbf{2{,}42\text{ atm}}\]

\textbf{Point 6} (adiabatique 5-6, $T_5=973{,}15$ K, jusqu'à $p_6=p_1=1$ atm) :
\[T_6=T_5\left(\frac{p_1}{p_4}\right)^{0{,}2857}=973{,}15\times\left(\frac{1}{2{,}42}\right)^{0{,}2857}
=973{,}15\times0{,}7769=\mathbf{756{,}0\text{ K}}\]

\textbf{Rendement.} Chaleur reçue (2-3 et 4-5, isobares) :
\[q_{23}=c_p(T_3-T_2)=1004\times591{,}93=594{,}3\text{ kJ/kg}\qquad
q_{45}=c_p(T_5-T_4)=1004\times93{,}07=93{,}4\text{ kJ/kg}\]
\[Q_{in}=q_{23}+q_{45}=\mathbf{687{,}7\text{ kJ/kg}}\]
Seul le travail de la 2\textsuperscript{nde} turbine est net (celui de la 1\textsuperscript{ère} est
entièrement consommé par le compresseur) :
\[|W_{net}|=c_p(T_5-T_6)=1004\times(973{,}15-756{,}0)=218{,}0\text{ kJ/kg}\]
\eqbox{\eta=\frac{|W_{net}|}{Q_{in}}=\frac{218{,}0}{687{,}7}=\mathbf{31{,}7\%}}

\textbf{Puissance utile de la 2\textsuperscript{nde} turbine} (convention : travail \emph{fourni par} le
fluide $\Rightarrow$ négatif) :
\eqbox{P=-\dot m\cdot|W_{net}|=-35\times218{,}0=\mathbf{-7{,}63\ MW}}
\end{correction}

\subsection{Exercice 2 -- Moteur diesel 4 cylindres}

\begin{enonce}
Moteur diesel 4 temps, 4 cylindres, cylindrée totale 2 L ($0{,}5$ L/cylindre). Rapport volumétrique
$\varepsilon=18$, rapport de combustion $\rho=2$. Air admis à $27\,^\circ$C, 1 bar, $\gamma=1{,}4$,
$R=8{,}314$ J/mol·K. Vitesse de rotation $N=2200$ tr/min.

\emph{Trouver $T,p,V$ en B, C, D, E ; le rendement théorique ; la chaleur $Q_{CD}$ et le travail théorique
par cylindre et par cycle ; la puissance théorique du moteur.}
\end{enonce}

\begin{correction}
\textbf{Volumes.} $V_B-V_A=500\text{ cm}^3$, $\varepsilon=V_B/V_A=18\Rightarrow V_A=V_C=29{,}4\text{ cm}^3$,
$V_B=529{,}4\text{ cm}^3$. Combustion isobare : $V_D=\rho V_C=58{,}8\text{ cm}^3$.

\textbf{B} : $T_B=300{,}15$ K, $p_B=1$ bar.

\textbf{C} (compression adiabatique) :
\[T_C=T_B\,\varepsilon^{\gamma-1}=300{,}15\times18^{0{,}4}=\mathbf{953{,}7\text{ K}}\qquad
p_C=p_B\,\varepsilon^\gamma=18^{1{,}4}=\mathbf{57{,}2\text{ bar}}\]

\textbf{D} (combustion isobare, $p_D=p_C$) : $T_D=T_C\cdot\rho=953{,}7\times2=\mathbf{1907{,}4\text{ K}}$

\textbf{E} (détente adiabatique jusqu'à $V_E=V_B$, soit $V_D/V_E=\rho/\varepsilon=1/9$) :
\[T_E=T_D\left(\frac{1}{9}\right)^{0{,}4}=\mathbf{792{,}2\text{ K}}\qquad
p_E=p_D\left(\frac{1}{9}\right)^{1{,}4}=\mathbf{2{,}64\text{ bar}}\]

\textbf{Rendement théorique} :
\[\eta_{\text{théo}}=1-\frac{1}{\varepsilon^{\gamma-1}}\cdot\frac{\rho^\gamma-1}{\gamma(\rho-1)}
=1-\frac{1}{3{,}178}\times\frac{2^{1{,}4}-1}{1{,}4\times1}=1-0{,}3147\times1{,}171\]
\eqbox{\eta_{\text{théo}}=\mathbf{63{,}2\%}}

\emph{Remarque : c'est exactement la formule du formulaire
$\eta=1+\dfrac{\varepsilon^{1-\gamma}-\varepsilon^{1-\gamma}\rho^{\gamma}}{\gamma(\rho-1)}$, simplement
réarrangée en mettant $\varepsilon^{1-\gamma}=\frac{1}{\varepsilon^{\gamma-1}}$ en facteur. Utilisez celle
que vous préférez, elles donnent le même nombre.}

\textbf{Chaleur apportée} (isobare, $n=\dfrac{p_BV_B}{RT_B}=\dfrac{10^5\times529{,}4\times10^{-6}}
{8{,}314\times300{,}15}=0{,}02123$ mol, $C_p=\dfrac{\gamma R}{\gamma-1}=29{,}1$ J/mol·K) :
\[Q_{CD}=nC_p(T_D-T_C)=0{,}02123\times29{,}1\times953{,}7=\mathbf{589\text{ J}}\ \text{(par cylindre, par cycle)}\]

\textbf{Travail théorique} : $|W_{\text{théo}}|=\eta_{\text{théo}}\cdot Q_{CD}=0{,}632\times589=\mathbf{372\text{ J}}$

\textbf{Puissance théorique} ($z=4$, 1 cycle = 2 tours) :
\eqbox{P_{th}=\frac{N}{2\cdot60}\cdot|W_{\text{théo}}|\cdot z=\frac{2200}{120}\times372\times4
=\mathbf{27{,}3\ kW}}
\end{correction}

\subsection{Exercice 3 -- Petit moteur 2 temps}

\begin{enonce}
Petit moteur 2 temps monocylindre, $\varepsilon=9$, $\gamma=1{,}4$. La combustion apporte $Q_{CD}=80$ J par
cycle. Rendement de forme $\eta_f=0{,}70$, rendement mécanique $\eta_{\text{méca}}=0{,}85$.
Vitesse de rotation $N=4000$ tr/min.

\emph{Trouver le rendement théorique, le travail disponible par cycle, et la puissance effective.}
\end{enonce}

\begin{correction}
\textbf{Rendement théorique} (même formule que l'essence 4 temps, seule la mécanique de renouvellement
du mélange change) :
\eqbox{\eta_{\text{théo}}=1-\varepsilon^{1-\gamma}=1-9^{-0{,}4}=1-0{,}4153=\mathbf{58{,}5\%}}

\textbf{Travaux en cascade} :
\[|W_{\text{théo}}|=\eta_{\text{théo}}\cdot Q_{CD}=0{,}585\times80=46{,}8\text{ J}\]
\[|W_{\text{indiqué}}|=\eta_f\cdot|W_{\text{théo}}|=0{,}70\times46{,}8=32{,}8\text{ J}\]
\[|W_{disponible}|=\eta_{\text{méca}}\cdot|W_{\text{indiqué}}|=0{,}85\times32{,}8=\mathbf{27{,}8\text{ J}}\]

\textbf{Puissance effective} -- \textbf{attention} : en 2 temps, 1 tour = 1 cycle, donc pas de
$\div 2$ :
\eqbox{P_{eff}=\frac{N}{60}\cdot|W_{disponible}|\cdot z=\frac{4000}{60}\times27{,}8\times1
=\mathbf{1{,}85\ kW}}
\end{correction}

\begin{piege}
C'est l'erreur la plus fréquente sur cette matière : garder le $\div 2\cdot 60$ du 4 temps pour un moteur
2 temps. Relisez toujours l'énoncé : « 1 tour = 1 cycle » (2 temps) ou « 1 cycle = 2 tours » (4 temps).
\end{piege}

\subsection{Exercice 4 -- Cycle frigorifique au NH\textsubscript{3}}

\begin{enonce}
Une machine frigorifique fonctionne au NH$_3$ (ammoniac). Évaporation à $-10\,^\circ$C, condensation à
$50\,^\circ$C, surchauffe et sous-refroidissement de $5\,^\circ$C chacun.

\emph{Placez le cycle sur le diagramme $(\log p,h)$ du NH$_3$ (fourni à l'examen), trouvez les enthalpies
de chaque point, puis la $PSN$.}
\end{enonce}

\begin{correction}
\textbf{Méthode} (voir section 2.2) : on repère $p_{\text{évap}}$ (isobare à $-10\,^\circ$C) et $p_{cond}$ (isobare
à $50\,^\circ$C) sur la cloche de saturation du NH$_3$.
\begin{itemize}[leftmargin=1.4em]
\item \textbf{1'} : sur l'isobare basse pression, à $-10+5=-5\,^\circ$C (zone vapeur, à droite de la
cloche).
\item \textbf{2'} : sur l'isobare haute pression, en suivant l'isentropique passant par 1'.
\item \textbf{3'} : sur l'isobare haute pression, à $50-5=45\,^\circ$C (zone liquide, à gauche de la cloche).
\item \textbf{4'} : à la verticale de 3' (même $h$), sur l'isobare basse pression.
\end{itemize}

\emph{Valeurs indicatives lues sur le diagramme (elles varient selon la précision de votre tracé -- l'ordre
de grandeur et la méthode comptent) :}
\[h_{1'}\approx1760\text{ kJ/kg}\qquad h_{2'}\approx2050\text{ kJ/kg}\qquad
h_{3'}=h_{4'}\approx700\text{ kJ/kg}\]

\textbf{PSN} :
\eqbox{PSN=\frac{Q_2}{W}=\frac{h_{1'}-h_{4'}}{h_{2'}-h_{1'}}=\frac{1760-700}{2050-1760}
=\frac{1060}{290}=\mathbf{3{,}66}}

\emph{Interprétation : pour 1 kWh de travail de compression fourni, on extrait $\approx3{,}66$ kWh à la
source froide -- c'est pour cela qu'une machine frigorifique/PAC est plus efficace qu'un chauffage
électrique direct (qui a un « rendement » de 1).}
\end{correction}

\begin{methode}
\textbf{Ce qu'il faut absolument savoir refaire sans le cours} pour les 3 exercices récurrents annoncés
(gaz réfrigérant type R134a, frigo/chauffage avec échange thermique, moteur essence/diesel/2 temps) :
\begin{enumerate}[leftmargin=1.6em]
\item Construire le tableau d'état et identifier chaque transformation (section 1).
\item Placer un cycle frigo sur un diagramme $(\log p,h)$ et en tirer $PSN$/$COP$ (section 2).
\item Dériver $\eta_{\text{théo}}$ à partir de $\varepsilon$ (et $\rho$) pour un moteur, \emph{sans} passer par
la combustion (section 1).
\item Relier un cycle fermé (cylindre, fluide frigorigène) à un circuit ouvert externe (eau de chauffage,
air ambiant) via la puissance échangée et $Q=\dot m\,c\,\Delta T$.
\end{enumerate}
\end{methode}
